\section{Results}
\lblsec{results}
We first compare our approach against recent methods for unpaired image-to-image translation on paired datasets where ground truth input-output pairs are available for evaluation. We then study the importance of both the adversarial loss and the cycle consistency loss and compare our full method against several variants. Finally, we demonstrate the generality of our algorithm on a wide range of applications where paired data does not exist. For brevity, we refer to our method as {\tt CycleGAN}. The \href{https://github.com/junyanz/pytorch-CycleGAN-and-pix2pix}{PyTorch} and \href{https://github.com/junyanz/CycleGAN}{Torch} code, models, and full results can be found at our \href{https://junyanz.github.io/CycleGAN/}{website}.



\subsection{Evaluation}
\lblsec{comparison}

\begin{figure*}[t]
\begin{center}
\includegraphics[width=0.95\linewidth]{figs/baseline_cityscapes_withfeatureandl1.pdf}
\end{center}
 \vspace{-4 mm}
 \caption{Different methods for mapping labels$\leftrightarrow$photos trained on Cityscapes images. From left to right: input, BiGAN/ALI~\cite{donahue2016adversarial,dumoulin2016adversarially}, CoGAN~\cite{liu2016coupled},  feature loss + GAN, SimGAN~\cite{shrivastava2016learning}, CycleGAN (ours), pix2pix~\cite{isola2016image} trained on paired data, and ground truth.}
  \vspace{-4 mm}
\lblfig{cityscapes_baseline}
\end{figure*}




\begin{figure*}[t]
\begin{center}
\includegraphics[width=0.95\linewidth]{figs/maps_comparison_withfeatureandl1.pdf} 
 \vspace{-4 mm}
\caption{Different methods for mapping aerial photos$\leftrightarrow$maps on Google Maps. From left to right: input, BiGAN/ALI~\cite{donahue2016adversarial,dumoulin2016adversarially}, CoGAN~\cite{liu2016coupled},  feature loss + GAN, SimGAN~\cite{shrivastava2016learning}, CycleGAN (ours), pix2pix~\cite{isola2016image} trained on paired data, and ground truth.}
\lblfig{maps_baseline}
\end{center}
 \vspace{-8 mm}

\end{figure*} 


Using the same evaluation datasets and metrics as ``pix2pix''~\shortcite{isola2016image}, we compare our method against several baselines both qualitatively and quantitatively. The tasks include semantic labels$\leftrightarrow$photo on the Cityscapes dataset~\cite{Cordts2016Cityscapes}, and map$\leftrightarrow$aerial photo on data scraped from Google Maps. We also perform ablation study on the full loss function.
	
\vspace{-1 mm}
\subsubsection{Evaluation Metrics}
\vspace{-1 mm}

\hspace{4mm} {\bf AMT perceptual studies} On the map$\leftrightarrow$aerial photo task, we run ``real vs fake" perceptual studies on Amazon Mechanical Turk (AMT) to assess the realism of our outputs. We follow the same perceptual study protocol from Isola et al.~\shortcite{isola2016image}, except we only gather data from $25$ participants per algorithm we tested. Participants were shown a sequence of pairs of images, one a real photo or map and one fake (generated by our algorithm or a baseline), and asked to click on the image they thought was real. The first $10$ trials of each session were practice and feedback was given as to whether the participant's response was correct or incorrect. The remaining $40$ trials were used to assess the rate at which each algorithm fooled participants. Each session only tested a single algorithm, and participants were only allowed to complete a single session. The numbers we report here are not directly comparable to those in~\cite{isola2016image} as our ground truth images were processed slightly differently
\footnote{We train all the models on $256\times256$ images while in pix2pix~\cite{isola2016image}, the model was trained on $256\times256$ patches of $512\times512$ images, and run convolutionally on the $512\times512$ images at test time. We choose $256\times256$ in our experiments as many baselines cannot scale up to high-resolution images, and CoGAN cannot be tested fully convolutionally.} and the participant pool we tested may be differently distributed from those tested in~\cite{isola2016image} (due to running the experiment at a different date and time). Therefore, our numbers should only be used to compare our current method against the baselines (which were run under identical conditions), rather than against~\cite{isola2016image}. 

{\bf FCN score} Although perceptual studies may be the gold standard for assessing graphical realism, we also seek an automatic quantitative measure that does not require human experiments. For this, we adopt the ``FCN score" from~\cite{isola2016image}, and use it to evaluate the Cityscapes labels$\rightarrow$photo task. The FCN metric evaluates how interpretable the generated photos are according to an off-the-shelf semantic segmentation algorithm (the fully-convolutional network, FCN, from~\cite{long2015fully}). The FCN predicts a label map for a generated photo. This label map can then be compared against the input ground truth labels using standard semantic segmentation metrics described below. The intuition is that if we generate a photo from a label map of ``car on the road", then we have succeeded if the FCN applied to the generated photo detects ``car on the road".

{\bf Semantic segmentation metrics} To evaluate the performance of photo$\rightarrow$labels, we use the standard metrics from the Cityscapes benchmark~\cite{Cordts2016Cityscapes}, including per-pixel accuracy, per-class accuracy, and mean class Intersection-Over-Union (Class IOU)~\cite{Cordts2016Cityscapes}.

\vspace{-1 mm}
\subsubsection{Baselines}
\vspace{-1 mm}

\hspace{4mm} {\bf CoGAN~\cite{liu2016coupled}} This method learns one GAN generator for domain $X$ and one for domain $Y$, with tied weights on the first few layers for shared latent representations. Translation from $X$ to $Y$ can be achieved by finding a latent representation that generates image $X$ and then rendering this latent representation into style $Y$.


{\bf SimGAN~\cite{shrivastava2016learning}} Like our method, Shrivastava et al.\shortcite{shrivastava2016learning} uses an adversarial loss to train a translation from $X$ to $Y$. The regularization term $\norm{x-G(x)}_1$ i s used to penalize making large changes at pixel level.

{\bf Feature loss + GAN} We also test a variant of SimGAN~\cite{shrivastava2016learning} where the L1 loss is computed over deep image features using a pretrained network (VGG-16 \texttt{relu4\_2}~\cite{simonyan2014very}), rather than over RGB pixel values. Computing distances in deep feature space, like this, is also sometimes referred to as using a ``perceptual loss"~\cite{dosovitskiy2016generating,johnson2016perceptual}.

{\bf BiGAN/ALI~\cite{dumoulin2016adversarially,donahue2016adversarial}} Unconditional GANs~\cite{goodfellow2014generative} learn a generator $G: Z\rightarrow X$, that maps a random noise $z$ to an image $x$. The BiGAN~\cite{dumoulin2016adversarially} and ALI~\cite{donahue2016adversarial} propose to also learn the inverse mapping function $F: X\rightarrow Z$. Though they were originally designed for mapping a latent vector $z$ to an image $x$, we implemented the same objective for mapping a source image $x$ to a target image $y$.

{\bf pix2pix~\cite{isola2016image}} We also compare against pix2pix~\cite{isola2016image}, which is trained on paired data, to see how close we can get to this ``upper bound" without using any paired data. 

For a fair comparison, we implement all the baselines using the same architecture and details as our method, except for CoGAN~\cite{liu2016coupled}. CoGAN builds on generators that produce images from a shared latent representation, which is incompatible with our image-to-image network. We use the public \href{https://github.com/mingyuliutw/CoGAN}{implementation} of CoGAN instead.



\begin{table}
\centering
\scalebox{0.75} {
\begin{tabular}{lcc}
 & \textbf{Map $\rightarrow$ Photo}  &  \textbf{Photo $\rightarrow$ Map} \\ 
\textbf{Loss} & \% Turkers labeled \emph{real} & \% Turkers labeled \emph{real} \\
\hline
CoGAN~\cite{liu2016coupled} & 0.6\% $\pm$ 0.5\% & 0.9\% $\pm$ 0.5\% \\
BiGAN/ALI~\cite{dumoulin2016adversarially,donahue2016adversarial} & 2.1\% $\pm$ 1.0\% & 1.9\% $\pm$ 0.9\% \\
SimGAN~\cite{shrivastava2016learning} & 0.7\% $\pm$ 0.5\% & 2.6\% $\pm$ 1.1\% \\
Feature loss + GAN & 1.2\% $\pm$ 0.6\% & 0.3\% $\pm$ 0.2\% \\
CycleGAN (ours)  & {\bf 26.8\% $\pm$ 2.8\%} &  {\bf 23.2\% $\pm$ 3.4\%} \\
\end{tabular} }
\vspace{-0.1in}
\caption {AMT ``real vs fake" test on maps$\leftrightarrow$aerial photos at $256\times 256$ resolution.}
\vspace{-0.1in}
\lbltbl{AMT_map2sat}
\end{table}


\begin{table}
\centering
\vspace{-2mm}
\scalebox{0.75} {
\begin{tabular}{lcccc}
 & & & \\
\textbf{Loss} & \textbf{Per-pixel acc.} & \textbf{Per-class acc.} & \textbf{Class IOU} \\ \hline
CoGAN~\cite{liu2016coupled} & 0.40 & 0.10 & 0.06 \\
BiGAN/ALI~\cite{dumoulin2016adversarially,donahue2016adversarial} & 0.19 & 0.06 & 0.02 \\
SimGAN~\cite{shrivastava2016learning} & 0.20 & 0.10 & 0.04 \\
Feature loss + GAN & 0.06 & 0.04 & 0.01 \\
CycleGAN (ours) & {\bf 0.52} & {\bf 0.17} & {\bf 0.11} \\ \hline
pix2pix~\cite{isola2016image} & 0.71 & 0.25 & 0.18 \\ 
\end{tabular} }
\vspace{-0.1in}
\caption {FCN-scores for different methods, evaluated on Cityscapes labels$\rightarrow$photo.}
\lbltbl{labels_to_image_results}
\vspace{-0.2in}
\end{table}


\begin{table}
\centering
\scalebox{0.75} {
\begin{tabular}{lccccc}
 & & & \\
\textbf{Loss} & \textbf{Per-pixel acc.} & \textbf{Per-class acc.} & \textbf{Class IOU} \\ \hline
CoGAN~\cite{liu2016coupled} & 0.45 & 0.11 & 0.08\\
BiGAN/ALI~\cite{dumoulin2016adversarially,donahue2016adversarial} & 0.41 & 0.13 & 0.07 \\
SimGAN~\cite{shrivastava2016learning} & 0.47 & 0.11 & 0.07\\
Feature loss + GAN & 0.50  & 0.10 & 0.06 \\
CycleGAN (ours) & {\bf 0.58} & {\bf 0.22} & {\bf 0.16}\\\hline
pix2pix~\cite{isola2016image} & 0.85 & 0.40 & 0.32 \\
\end{tabular} }
\vspace{-0.1in}
\caption {Classification performance of photo$\rightarrow$labels for different methods on cityscapes.}
\vspace{-0.1in}
\lbltbl{image_to_labels_results}
\end{table}


\begin{figure*}[ht]
\begin{center}
\includegraphics[width=1.0\linewidth]{figs/ablation_study.pdf}
\end{center}
 \vspace{-5 mm}
 \caption{Different variants of our method for mapping labels$\leftrightarrow$photos trained on cityscapes. From left to right: input, cycle-consistency loss alone, adversarial loss alone, GAN + forward cycle-consistency loss ($F(G(x)) \approx x$), GAN + backward cycle-consistency loss ($G(F(y)) \approx y$), CycleGAN (our full method), and ground truth. Both  \emph{Cycle alone} and \emph{GAN + backward} fail to produce images similar to the target domain. \emph{GAN alone} and \emph{GAN + forward} suffer from mode collapse, producing identical label maps regardless of the input photo.} 
\lblfig{ablation} 
 \vspace{-5 mm}
\end{figure*}








\begin{table}
\centering
\scalebox{0.75} {
\begin{tabular}{lcccc}
 & & & \\
\textbf{Loss} & \textbf{Per-pixel acc.} & \textbf{Per-class acc.} & \textbf{Class IOU} \\ \hline
Cycle alone & 0.22 & 0.07 & 0.02 \\
GAN alone & 0.51 & 0.11 & 0.08 \\
GAN + forward cycle & {\bf 0.55} & {\bf 0.18} & {\bf 0.12} \\ 
GAN + backward cycle & 0.39 & 0.14 & 0.06 \\ 
CycleGAN (ours) & 0.52 & 0.17 & 0.11 \\  
\end{tabular} }
\vspace{-0.1in}
\caption {Ablation study: FCN-scores for different variants of our method, evaluated on Cityscapes labels$\rightarrow$photo.}
\vspace{-0.2in}
\lbltbl{loss_fcn}
\end{table}



\begin{table}
\centering
\scalebox{0.75} {
\begin{tabular}{lcccc}
 & & & \\
\textbf{Loss} & \textbf{Per-pixel acc.} & \textbf{Per-class acc.} & \textbf{Class IOU} \\ \hline
Cycle alone & 0.10 & 0.05 & 0.02 \\
GAN alone & 0.53 & 0.11 & 0.07 \\
GAN + forward cycle & 0.49 & 0.11 & 0.07 \\ 
GAN + backward cycle & 0.01 & 0.06 & 0.01 \\ 
CycleGAN (ours)& {\bf 0.58} & {\bf 0.22} & {\bf 0.16} \\  
\end{tabular} }
\vspace{-0.1in}
\caption {Ablation study: classification performance of photo$\rightarrow$labels for different losses, evaluated on Cityscapes. }
\vspace{-0.2in}
\lbltbl{loss_clf}
\end{table}

\subsubsection{Comparison against baselines}
As can be seen in \reffig{cityscapes_baseline} and \reffig{maps_baseline}, we were unable to achieve compelling results with any of the baselines. Our method, on the other hand, can produce translations that are often of similar quality to the fully supervised pix2pix.

\reftbl{AMT_map2sat} reports performance regarding the AMT perceptual realism task. Here, we see that our method can fool participants on around a quarter of trials, in both the maps$\rightarrow$aerial photos direction and the aerial photos$\rightarrow$maps direction at $256\times 256$ resolution\footnote{We also train CycleGAN and pix2pix at $512\times 512$ resolution, and observe the comparable performance:  maps$\rightarrow$aerial photos: CycleGAN: $37.5\% \pm 3.6\%$ and pix2pix: $33.9\% \pm 3.1\%$; aerial photos$\rightarrow$maps: CycleGAN: $16.5\% \pm 4.1\%$ and pix2pix: $8.5\% \pm 2.6\%$}. All the baselines almost never fooled participants.

\reftbl{labels_to_image_results} assesses the performance of the labels$\rightarrow$photo task on the Cityscapes and \reftbl{image_to_labels_results} evaluates the opposite mapping (photos$\rightarrow$labels). In both cases, our method again outperforms the baselines.

\vspace{-1 mm}
\subsubsection{Analysis of the loss function}
\lblsec{ablation}
\vspace{-1 mm}


In \reftbl{loss_fcn} and \reftbl{loss_clf}, we compare against ablations of our full loss. Removing the GAN loss substantially degrades results, as does removing the cycle-consistency loss. We therefore conclude that both terms are critical to our results. We also evaluate our method with the cycle loss in only one direction: GAN + forward cycle loss $\mathbb{E}_{x\sim p_{\text{data}}(x)}[\norm{F(G(x))-x}_1]$, or GAN + backward cycle loss $\mathbb{E}_{y\sim p_{\text{data}}(y)}[\norm{G(F(y))-y}_1]$ (\refeq{cycle}) and find that it often incurs training instability and causes mode collapse, especially for the direction of the mapping that was removed. \reffig{ablation} shows several qualitative examples. 

\vspace{-1 mm}
\subsubsection{Image reconstruction quality}
\vspace{-1 mm}

In \reffig{reconstruction}, we show a few random samples of the reconstructed images $F(G(x))$. We observed that the reconstructed images were often close to the original inputs $x$, at both training and testing time, even in cases where one domain represents significantly more diverse information, such as map$\leftrightarrow$aerial photos. 


\begin{figure}[t]
\begin{center}
\includegraphics[width=1.0\linewidth]{figs/facades_and_sketches.pdf}
\end{center}
 \vspace{-5 mm}
 \caption{Example results of CycleGAN on paired datasets used in ``pix2pix''~\cite{isola2016image} such as architectural labels$\leftrightarrow$photos and edges$\leftrightarrow$shoes.}
\lblfig{facades_and_sketches}
 \vspace{-5 mm}
\end{figure}

\vspace{-1 mm}
\subsubsection{Additional results on paired datasets}
\vspace{-1 mm}
\reffig{facades_and_sketches} shows some example results on other paired datasets used in ``pix2pix''~\cite{isola2016image}, such as architectural labels$\leftrightarrow$photos from the CMP Facade Database~\cite{Tylecek13}, and edges$\leftrightarrow$shoes from the UT Zappos50K
dataset~\cite{yu2014fine}. The image quality of our results is close to those produced by the fully supervised pix2pix while our method learns the mapping without paired supervision. 


\vspace{-1 mm}
\subsection{Applications}
\lblsec{application}
\vspace{-1 mm}

We demonstrate our method on several applications where paired training data does not exist. Please refer to the appendix (\refsec{appendix}) for more details about the datasets. We observe that translations on training data are often more appealing than those on test data, and full results of all applications on both training and test data can be viewed on our project \href{https://junyanz.github.io/CycleGAN/}{website}.



{\bf Collection style transfer (\reffig{painting} and ~\reffig{painting2})}  We train the model on landscape photographs downloaded from Flickr and WikiArt. Unlike recent work on ``neural style transfer"~\cite{gatys2015neural}, our method learns to mimic the style of an entire \emph{collection} of artworks, rather than transferring the style of a single selected piece of art. Therefore, we can learn to generate photos in the style of, e.g., Van Gogh, rather than just in the style of Starry Night. The size of the dataset for each artist/style was $526$, $1073$, $400$, and $563$ for Cezanne, Monet, Van Gogh, and Ukiyo-e.

{\bf Object transfiguration (\reffig{big})} The model is trained to translate one object class from ImageNet~\cite{deng2009imagenet} to another (each class contains around $1000$ training images). Turmukhambetov et al.~\shortcite{turmukhambetov2015modeling} propose a subspace model to translate one object into another object of the same category, while our method focuses on object transfiguration between two visually similar categories. 

{\bf Season transfer (\reffig{big})}  The model is trained on $854$ winter photos and $1273$ summer photos of Yosemite downloaded from Flickr.

\begin{figure}[t]
\begin{center}
\includegraphics[width=1.0\linewidth]{figs/identity.pdf}
\end{center}
 \vspace{-5 mm}
 \caption{The effect of the \textit{identity mapping loss} on Monet's painting$\rightarrow$ photos. From left to right: input paintings, CycleGAN without identity mapping loss, CycleGAN with identity mapping loss. The identity mapping loss helps preserve the color of the input paintings. }
\lblfig{identity}
 \vspace{-6mm}
\end{figure}

{\bf Photo generation from paintings (\reffig{painting2photo})} 
For painting$\rightarrow$photo, we find that it is helpful to introduce an additional loss to encourage the mapping to preserve color composition between the input and output. In particular, we adopt the technique of Taigman et al.~\cite{taigman2016unsupervised} and regularize the generator to be near an identity mapping when real samples of the target domain are provided as the input to the generator: i.e., $ \mathcal{L}_{\text{identity}}(G, F) = \mathbb{E}_{y\sim p_{\text{data}}(y)}[\norm{G(y) - y}_1]+\mathbb{E}_{x\sim p_{\text{data}}(x)}
    [\norm{F(x) - x}_1].$ 

Without $\mathcal{L}_{\text{identity}}$, the generator $G$ and $F$ are free to change the tint of input images when there is no need to. For example, when learning the mapping between Monet's paintings and Flickr photographs, the generator often maps paintings of daytime to photographs taken during sunset, because such a mapping may be equally valid under the adversarial loss and cycle consistency loss. 
The effect of this \emph{identity mapping loss} are shown in \reffig{identity}.

In \reffig{painting2photo}, we show additional results translating Monet's paintings to photographs. This figure and \reffig{identity} show results on paintings that were included in the \emph{training set}, whereas for all other experiments in the paper, we only evaluate and show test set results. Because the training set does not include paired data, coming up with a plausible translation for a training set painting is a nontrivial task. Indeed, since Monet is no longer able to create new paintings, generalization to unseen, ``test set", paintings is not a pressing problem.

{\bf Photo enhancement (\reffig{flower})} We show that our method can be used to generate photos with shallower depth of field. We train the model on flower photos downloaded from Flickr. The source domain consists of flower photos taken by smartphones, which usually have deep DoF due to a small aperture. The target contains photos captured by DSLRs with a larger aperture. Our model successfully generates photos with shallower depth of field from the photos taken by smartphones. 


{\bf Comparison with Gatys et al.~\cite{gatys2015neural}}
In ~\reffig{gatys_style}, we compare our results with neural style transfer~\cite{gatys2015neural} on photo stylization. For each row, we first use two representative artworks as the style images for ~\cite{gatys2015neural}. Our method, on the other hand, can produce photos in the style of entire \emph{collection}. To compare against neural style transfer of an entire collection, we compute the average Gram Matrix across the target domain and use this matrix to transfer the ``average style" with Gatys et al~\cite{gatys2015neural}. 

\reffig{gatys_others} demonstrates similar comparisons for other translation tasks. 
We observe that Gatys et al.~\cite{gatys2015neural} requires finding target style images that closely match the desired output, but still often fails to produce photorealistic results, while our method succeeds to generate natural-looking results, similar to the target domain. 